% Source: https://artofproblemsolving.com/community/c3h3084312_my_3000th_post
% Stable identity: aops-topic-3084312-q03
A number with 12 divisors has prime-exponent pattern $(11),(5,1),(3,2)$, or $(2,1,1)$. The least values for the first three patterns are $2^{11}=2048$, $2^5\cdot3=96$, and $2^3\cdot3^2=72$. For the last pattern the smallest is $2^2\cdot3\cdot5=60$ and the next smallest is $2^2\cdot3\cdot7=84$; changing which prime is squared gives at least $2\cdot3^2\cdot5=90$, and using a different prime set gives at least 84. In pattern $(3,2)$, any other assignment gives at least $2^2\cdot3^3=108$. Thus the first three are 60,72,84, with sum $\boxed{216}$.
